\chapter{Sets, Relations, Functions, and Cardinality}

A security tester receives a model specification: a classifier with thirty binary
features, each taking only the value zero or one. To test exhaustively, it must try
every possible input. Each feature doubles the number of possibilities, so the answer
is $2^{30} = 1{,}073{,}741{,}824$. At one million inputs per second, that test takes
$1073.74$ seconds, about $17.9$ minutes. Add ten more features: the number becomes
$2^{40} = 1{,}099{,}511{,}627{,}776$, and a machine a thousand times faster needs
$1099.51$ seconds, about $18.3$ minutes. We multiply the problem by a thousand and
speed the machine up a thousandfold, and we are back at the starting point.

The same difficulty appears more sharply once features may take an arbitrary real
value. The number of inputs stops being "many" and starts being "infinitely many",
and the question changes: can two infinite sets still be compared in size? Instinct
says there are more rational numbers than natural numbers, since between any two
natural numbers there always lie infinitely many rationals. That instinct is wrong,
and this chapter shows why.

The previous chapter gave us a language: quantifiers, implication, proof patterns.
What was missing is the object being talked about. This chapter installs its basic
vocabulary, namely sets and their operations, relations, functions, and one way of
comparing the size of sets that does not depend on the ability to count. That
vocabulary is used throughout the book, from subspaces in Chapter 6 to formal proofs
in Chapter 28.

\begin{untukanak}
Picture two baskets of marbles, one red and one blue. We want to know which basket
holds more, but we are not allowed to count. Take one red marble, pair it with one
blue marble, set both aside, and repeat until one basket is empty: if both empty at
the same moment they hold the same number, if blue is left over then blue holds more.
This method never mentions a number; it only pairs.

Now suppose the baskets never end, so that we can never finish moving marbles. Can
pairing still be used? It can, and that is the surprise of this chapter: with the
right pairing rule, all the counting numbers can be paired off with all the even
numbers, with nothing left over on either side, and pairs of numbers can be paired
with a single number. We will understand why pairing is stronger than counting, and
why some endless sets still cannot be paired with the counting numbers, however hard
we try.
\end{untukanak}

\section{Operations on sets and De Morgan's laws}
\tingkat{basic}

A single set is of no use until we can build new sets from it. Four ways of building
are defined by the same pattern: state a membership condition, and the new set is the
collection of all objects satisfying it. Let $X$ be a universal set and
$A, B \subseteq X$; from Chapter 1 we recall that $A \subseteq B$ means every member
of $A$ is also a member of $B$, and that equality of sets is proved by two inclusions.

\begin{definisi}[Union, intersection, difference, complement]
The union $A \cup B$ is the set of all $x \in X$ that are members of $A$ or members
of $B$. The intersection $A \cap B$ is the set of all $x \in X$ that are members of
$A$ and at the same time members of $B$. The difference $A \setminus B$ is the set of
all $x \in X$ that are members of $A$ but not members of $B$. The complement
$A^{c} = X \setminus A$ is the set of all $x \in X$ that are not members of $A$.
\end{definisi}

The word "or" in a union is inclusive: $x$ may lie in both sets, and that is the
condition under which $A \cap B \subseteq A \cup B$ is always true, as Hoffman and
Kunze note in the Appendix. The empty set $\varnothing$ is needed so that the
intersection of two disjoint sets is still a set.

\begin{proposisi}[De Morgan's laws]\index{De Morgan's laws}
For any $A, B \subseteq X$,
\[
  X \setminus (A \cup B) = (X \setminus A) \cap (X \setminus B),
  \qquad
  X \setminus (A \cap B) = (X \setminus A) \cup (X \setminus B).
\]
\end{proposisi}
\par\smallskip\noindent\textit{Proof.} Take any $x \in X$. By the definition of
complement and union,
\[
  x \in X \setminus (A \cup B)
  \iff x \in X \ \text{and}\ \text{not}\ (x \in A \ \text{or}\ x \in B).
\]
By De Morgan's law for statements from Chapter 1, the negation of "at least one" is
"both fail", so
\[
  x \in X \setminus (A \cup B)
  \iff x \in X \ \text{and}\ x \notin A \ \text{and}\ x \notin B
  \iff x \in (X \setminus A) \cap (X \setminus B).
\]
Since $x$ was arbitrary, the two sets are equal. The second equality is obtained by
replacing every "and" by "or" in the same chain.
\hfill$\blacksquare${\parfillskip=0pt\par}

The law holds for two sets; what we need later is the version for any number of sets.

\begin{definisi}[Indexed family]
Let $I$ be an index set and let $A_i \subseteq X$ be given for each $i \in I$.
\begin{align*}
  \bigcup_{i \in I} A_i &= \{\, x \in X : x \in A_i \ \text{for some}\ i \in I \,\}, \\
  \bigcap_{i \in I} A_i &= \{\, x \in X : x \in A_i \ \text{for every}\ i \in I \,\}.
\end{align*}
\end{definisi}

The general version reads $X \setminus \bigcup_{i} A_i = \bigcap_{i} (X \setminus A_i)$
and $X \setminus \bigcap_{i} A_i = \bigcup_{i} (X \setminus A_i)$. Its proof adds no
new idea: only "for some" turns into "for every" when the negation is pulled inside,
exactly as in Chapter 1.

Numbers make the definition checkable. For $X = \{1,\dots,10\}$,
$A = \{1,\dots,5\}$, and $B = \{4,\dots,8\}$, Figure~\ref{fig:venn} shows the three
operations on the same sets: the union has $3 + 2 + 3 = 8$ members, the intersection
has $2$, and the complement of $B$ has $3 + 2 = 5$.

\begin{figure}[htbp]
\centering
\begin{tikzpicture}[scale=0.82,every node/.style={font=\footnotesize}]
  \begin{scope}
    \begin{scope}
      \clip (-0.5,0) circle (0.85);
      \fill[layerteal!30] (-1.55,-1.25) rectangle (1.55,1.25);
    \end{scope}
    \begin{scope}
      \clip (0.5,0) circle (0.85);
      \fill[layerteal!30] (-1.55,-1.25) rectangle (1.55,1.25);
    \end{scope}
    \draw[ink,thick] (-0.5,0) circle (0.85);
    \draw[ink,thick] (0.5,0) circle (0.85);
    \draw[ink,thick] (-1.55,-1.25) rectangle (1.55,1.25);
    \node at (-1.0,0) {3};
    \node at (0,0) {2};
    \node at (1.0,0) {3};
    \node at (0,-1.0) {2};
    \node[above] at (-0.5,0.86) {$A$};
    \node[above] at (0.5,0.86) {$B$};
    \node[below] at (0,-1.25) {$A \cup B$: 8};
  \end{scope}
  \begin{scope}[xshift=3.9cm]
    \begin{scope}
      \clip (-0.5,0) circle (0.85);
      \clip (0.5,0) circle (0.85);
      \fill[accent!35] (-1.55,-1.25) rectangle (1.55,1.25);
    \end{scope}
    \draw[ink,thick] (-0.5,0) circle (0.85);
    \draw[ink,thick] (0.5,0) circle (0.85);
    \draw[ink,thick] (-1.55,-1.25) rectangle (1.55,1.25);
    \node at (-1.0,0) {3};
    \node at (0,0) {2};
    \node at (1.0,0) {3};
    \node at (0,-1.0) {2};
    \node[above] at (-0.5,0.86) {$A$};
    \node[above] at (0.5,0.86) {$B$};
    \node[below] at (0,-1.25) {$A \cap B$: 2};
  \end{scope}
  \begin{scope}[xshift=7.8cm]
    \fill[layergrey!22] (-1.55,-1.25) rectangle (1.55,1.25);
    \fill[white] (0.5,0) circle (0.85);
    \draw[ink,thick] (-0.5,0) circle (0.85);
    \draw[ink,thick] (0.5,0) circle (0.85);
    \draw[ink,thick] (-1.55,-1.25) rectangle (1.55,1.25);
    \node at (-1.0,0) {3};
    \node[ink!40] at (0,0) {2};
    \node[ink!40] at (1.0,0) {3};
    \node at (0,-1.0) {2};
    \node[above] at (-0.5,0.86) {$A$};
    \node[above] at (0.5,0.86) {$B$};
    \node[below] at (0,-1.25) {$B^{c}$: 5};
  \end{scope}
\end{tikzpicture}
\caption{Three Venn diagrams over the same sets, with the number of members in each
region. Left: $A \cup B$ has $3 + 2 + 3 = 8$ members; middle: $A \cap B$ has $2$
members; right: $B^{c}$ has $3 + 2 = 5$ members, namely the shaded region, while the
members inside $B$ are printed faded.}
\label{fig:venn}
\end{figure}

\begin{tangan}
Materials: twenty used name cards, two pieces of string two meters long, an open floor.

Form two string loops overlapping by about a quarter of their width. Write the names
of five friends who like coffee and five friends who like tea. Put the coffee cards
in the first loop, the tea cards in the second, then move the card of a friend who
likes both into the overlap: that is the intersection. Pick one name who likes
neither and put that card outside both loops.

Count three numbers: cards inside both loops, cards in the overlap, and cards outside
the second loop. What must become visible: the number of cards in the union is not the
number of coffee cards plus the number of tea cards, because the cards in the overlap
are counted twice, and the excess is exactly the intersection.
\end{tangan}

The set we most often build from $A$ is the set of all its subsets, and its size grows
far faster than $A$ itself.

\begin{definisi}[Power set]\index{power set}
The power set $\mathcal{P}(A)$ is the set of all subsets of $A$, that is,
$\mathcal{P}(A) = \{\, B : B \subseteq A \,\}$.
\end{definisi}

\kamus{power set}{himpunan kuasa}{A list of every way of choosing part of the members of a bag, including choosing nothing and choosing everything.}

For example, $\mathcal{P}(\varnothing) = \{\varnothing\}$ has one member, and
$\mathcal{P}(\{1,2\}) = \{\varnothing, \{1\}, \{2\}, \{1,2\}\}$ has four.

\begin{proposisi}
If $A$ is finite with $|A| = n$, then $|\mathcal{P}(A)| = 2^{n}$.
\end{proposisi}
\par\smallskip\noindent\textit{Proof.} For $B \subseteq A$ define the characteristic
function $\chi_B : A \to \{0,1\}$ by $\chi_B(a) = 1$ if $a \in B$ and
$\chi_B(a) = 0$ if $a \notin B$. The map $B \mapsto \chi_B$ is a bijection between
$\mathcal{P}(A)$ and the set of all functions $A \to \{0,1\}$: from $\chi$ we recover
$B = \{a : \chi(a) = 1\}$, and two different subsets give two different functions.
Hence $|\mathcal{P}(A)|$ equals the number of functions $A \to \{0,1\}$.

We count those functions by induction on $n$. For $n = 0$ there is only the empty
function, and $2^{0} = 1$. Suppose $|A| = n+1$, choose $a_0 \in A$, and write
$A' = A \setminus \{a_0\}$. Every $f : A \to \{0,1\}$ is determined uniquely by the
restriction $f|_{A'}$ and the value $f(a_0)$; the induction hypothesis gives $2^{n}$
choices of restriction and $2$ choices of value, $2^{n+1}$ in all.
\hfill$\blacksquare${\parfillskip=0pt\par}

That proof never mentions combinations: the bijection moves the count, and induction
splits it into two free choices.

\begin{bacakeras}
Read in pairs, one voice for member and one voice for subset. Start slowly, then
speed up.

Member: One orange goes into the basket. I am inside.
Subset: One basket goes into the store. I am inside.
Member: I am cargo.
Subset: I am container.
Member: I am pointed at one by one.
Subset: I am proved all at once.
Member: When I go in, I am a single mark.
Subset: When I go in, I am a single proof.
Member: I am small and many.
Subset: I am large and sometimes only one.

Repeat until the last line comes out without thinking. The line repeated three times
together: cargo goes into bag, bag goes into bag, and the two use different marks.
\end{bacakeras}

\section{Cartesian products and relations}
\tingkat{basic}

Order matters. The set $\{1,2\}$ equals $\{2,1\}$, but the coordinates of a point in
the plane may not be swapped without moving the point. We therefore need an object
that remembers order, and its definition must be built from sets alone.

\begin{definisi}[Ordered pair and Cartesian product]\index{Cartesian product}
The ordered pair $(a,b)$ is defined as $\{\{a\}, \{a,b\}\}$. The Cartesian product
$A \times B$ is the set of all ordered pairs $(a,b)$ with $a \in A$ and $b \in B$.
\end{definisi}

The definition looks strange, but only one property is needed, namely
$(a,b) = (c,d)$ if and only if $a = c$ and $b = d$. Checking that property requires
four cases, and we leave it as Exercise 4.

The size of a product is easy to guess, and its proof uses a cut along the first
coordinate.

\begin{proposisi}
If $A$ and $B$ are finite, then $|A \times B| = |A| \cdot |B|$.
\end{proposisi}
\par\smallskip\noindent\textit{Proof.} For each $a \in A$ the set
$\{a\} \times B$ has $|B|$ members, since $b \mapsto (a,b)$ is a bijection between $B$
and $\{a\} \times B$. The family $\{\{a\} \times B\}_{a \in A}$ consists of mutually
disjoint sets, because $(a,b) = (a',b')$ forces $a = a'$. Their union is $A \times B$,
so $|A \times B|$ is the sum of $|A|$ terms each equal to $|B|$.
\hfill$\blacksquare${\parfillskip=0pt\par}

Now the object that gives this chapter its title. A relation is merely a statement
that may be true or false for each pair of objects, and so it can be stored as the set
of pairs for which the statement is true.

\begin{definisi}[Relation]\index{relation}
A relation from $A$ to $B$ is any $R \subseteq A \times B$; when $A = B$ we call it a
relation on $A$. We write $a \mathrel{R} b$ for $(a,b) \in R$. The domain of $R$ is
$\{\, a \in A : (a,b) \in R\ \text{for some}\ b \in B \,\}$ and its range is
$\{\, b \in B : (a,b) \in R\ \text{for some}\ a \in A \,\}$.
\end{definisi}

Since a relation on $A$ is a subset of $A \times A$, the number of relations on a set
with $n$ members is $2^{n^{2}}$. For $n = 3$ the number is $2^{9} = 512$ and for
$n = 4$ it becomes $2^{16} = 65{,}536$; both are checked by the script in the lab
section.

Four properties filter them into useful objects. A relation $R$ on $A$ is called
reflexive if $a \mathrel{R} a$ for every $a \in A$; symmetric if
$a \mathrel{R} b$ implies $b \mathrel{R} a$; transitive if $a \mathrel{R} b$ and
$b \mathrel{R} c$ imply $a \mathrel{R} c$; and antisymmetric if
$a \mathrel{R} b$ together with $b \mathrel{R} a$ is possible only when $a = b$.

\begin{definisi}[Partial order and equivalence relation]\index{equivalence relation}
A relation on $A$ that is reflexive, antisymmetric, and transitive is called a partial
order. A relation on $A$ that is reflexive, symmetric, and transitive is called an
equivalence relation.
\end{definisi}

The relation $\le$ on $\mathbb{R}$ is a partial order, as is the relation "divides" on
the positive natural numbers, and $\subseteq$ on $\mathcal{P}(A)$ gives an example
that will later become the framework of subspaces in Chapter 6. The word "partial" is
used because two members need not be comparable: on $\mathcal{P}(\{1,2,3\})$ the sets
$\{1\}$ and $\{2\}$ are not contained in one another. An order in which every two
members are comparable is called a total order.

\section{Equivalence relations and partitions}
\tingkat{intermediate}

An equivalence relation cuts a set into groups that do not overlap. Before reaching
the theorem, we need a name for the groups.

\begin{definisi}[Equivalence class]\index{equivalence class}
Let $\sim$ be an equivalence relation on $A$ and $x \in A$. The equivalence class of
$x$ is $[x] = \{\, y \in A : y \sim x \,\}$.
\end{definisi}

Since $x \sim x$, no class is empty. The more interesting property: two classes cannot
overlap even slightly.

\begin{lemma}
Let $\sim$ be an equivalence relation on $A$ and $x, y \in A$. Then (i) $x \in [x]$;
(ii) $[x] = [y]$ if and only if $x \sim y$; and (iii) if
$[x] \cap [y] \neq \varnothing$, then $[x] = [y]$.
\end{lemma}
\par\smallskip\noindent\textit{Proof.} Part (i) is the statement $x \sim x$, true
because the relation is reflexive. For (ii), suppose $[x] = [y]$. By (i),
$y \in [y] = [x]$, so $y \sim x$ and by symmetry $x \sim y$. Conversely, suppose
$x \sim y$ and take $z \in [x]$, meaning $z \sim x$. From $z \sim x$ and
$x \sim y$ we get $z \sim y$, so $z \in [y]$; this gives
$[x] \subseteq [y]$, and the argument with the roles exchanged gives
$[y] \subseteq [x]$. For (iii), suppose $z \in [x] \cap [y]$. Then $z \sim x$ and
$z \sim y$; from $z \sim x$ and symmetry we get $x \sim z$, and together with
$z \sim y$ and transitivity we get $x \sim y$, so $[x] = [y]$ by (ii).
\hfill$\blacksquare${\parfillskip=0pt\par}

Consequently the set of all classes divides $A$ into pieces that are mutually disjoint
and cover all of $A$.

\begin{definisi}[Partition]\index{partition}
A partition of a set $A$ is a family $\mathcal{B}$ of nonempty subsets of $A$ such
that every $x \in A$ lies in exactly one member of $\mathcal{B}$.
\end{definisi}

\begin{teorema}
Every equivalence relation on $A$ determines a partition of $A$, every partition of
$A$ determines an equivalence relation on $A$, and the two processes invert each
other.
\end{teorema}
\par\smallskip\noindent\textit{Proof.} Let $\sim$ be an equivalence relation and
$\mathcal{B} = \{\, [x] : x \in A \,\}$. By the Lemma, each $[x]$ is nonempty and two
of its members are either disjoint or identical; every $x$ lies in $[x]$. Hence
$\mathcal{B}$ is a partition.

Conversely, let $\mathcal{B}$ be a partition of $A$ and set $x \sim y$ when $x$ and
$y$ lie in the same member of $\mathcal{B}$. Reflexivity and symmetry are clear. For
transitivity, suppose $x \sim y$ and $y \sim z$, and let $B$ be the member of
$\mathcal{B}$ containing $y$; then $x \in B$ and $z \in B$, so $x \sim z$.

The two processes invert each other: from $\sim$ we obtain the partition into classes,
and that partition returns the original $\sim$ because $x \sim y$ exactly when
$x \in [y]$; from a partition we obtain $\sim$, and the class $[x]$ is the member of
$\mathcal{B}$ containing $x$.
\hfill$\blacksquare${\parfillskip=0pt\par}

The number of classes is called the index. The set of all classes is written
$A/{\sim}$ and called the quotient set, while $q : A \to A/{\sim}$ with $q(x) = [x]$
is called the quotient map. Hoffman and Kunze note that every equivalence relation can
be viewed as coming from a function: two members are equivalent exactly when that
function gives them the same value.

\begin{contoh}
On the integers, set $x \equiv y \pmod{n}$ when $n$ divides $x-y$. This relation is
reflexive because $n \mid 0$, symmetric because $n \mid (x-y)$ implies
$n \mid (y-x)$, and transitive because $n \mid (x-y)$ together with
$n \mid (y-z)$ implies $n \mid (x-z)$. Its classes are the sets of numbers with the
same remainder, and there are exactly $n$ classes, namely $[0], \dots, [n-1]$, since
every integer can be written $x = qn + r$ with $0 \le r \le n-1$ uniquely.
\end{contoh}

An example closer to our own work: for $f : A \to B$, set $x \sim y$ when
$f(x) = f(y)$. The class $[x]$ is the set of inputs with the same output; in machine
learning that is "one prediction file", and the number of classes equals the number of
distinct labels.

\begin{untukanak}
The wardrobe at home is usually not sorted by any official rule. White shirts pile up
on one shelf, printed shirts on another, sports shirts in a third basket. Nobody
writes the rule down, yet every shirt sits in exactly one place.

That is what an equivalence relation does: it decides when two items count as
roommates, then groups all the items. The result is a tidy division, every item lands
in exactly one group, and different groups share no members.

Now the question that makes one curious. What if one shirt belongs to the white group
under one rule, but to the sports group under another rule? Is the grouping still
tidy? The tidiness does not depend on the rule chosen, but on three properties: every
item is a roommate of itself, roommatehood works both ways, and roommatehood can be
passed along through a third item.
\end{untukanak}

\section{Functions: injective, surjective, bijective}
\tingkat{intermediate}

We now fix the meaning of the word "function", because this chapter depends on the
fine distinction between codomain and range.

\begin{definisi}[Function]\index{injective function}
A function from $A$ to $B$ is a rule $f$ that pairs every $x \in A$ with exactly one
$f(x) \in B$. The set $A$ is called the domain, $B$ is called the codomain, and
$f(A) = \{\, f(x) : x \in A \,\}$ is called the range.
\end{definisi}

Codomain and range are different, and the difference matters. The function
$f : \mathbb{R} \to \mathbb{R}$ with $f(x) = x^{2}$ has codomain $\mathbb{R}$ but
range $[0,\infty)$. Hoffman and Kunze use three names interchangeably for the same
object: function, mapping, transformation.

For $S \subseteq B$ we define the preimage
$f^{-1}(S) = \{\, x \in A : f(x) \in S \,\}$. The notation $f^{-1}$ here does not
presuppose that $f$ has an inverse: a preimage is always defined, an inverse need not
exist.

\begin{definisi}[Injective, surjective, bijective]\index{surjective function}\index{bijection}
A function $f : A \to B$ is called injective if $f(x) = f(y)$ is possible only when
$x = y$; surjective if $f(A) = B$; and bijective if both.
\end{definisi}

\kamus{bijection}{bijeksi}{A perfect matching: every item in the first box has exactly one partner in the second box, with nothing left over on either side.}

Figure~\ref{fig:bijeksi} compares the three on two columns of nodes: the left panel
leaves no node over on either side, the middle panel leaves $d$ over on the right, and
the right panel sends two left nodes to one right node.

\begin{figure}[htbp]
\centering
\begin{tikzpicture}[scale=0.95,every node/.style={font=\footnotesize},>=Stealth,
  kotak/.style={draw=ink,thick,circle,inner sep=1.2pt,minimum size=4.4mm}]
  \begin{scope}
    \foreach \i/\y in {1/1.0,2/0,3/-1.0} {
      \node[kotak] (a\i) at (0,\y) {\i};
    }
    \node[kotak] (ba) at (1.9,1.0) {$a$};
    \node[kotak] (bb) at (1.9,0) {$b$};
    \node[kotak] (bc) at (1.9,-1.0) {$c$};
    \draw[->,thick,layerjade] (a1) -- (ba);
    \draw[->,thick,layerjade] (a2) -- (bb);
    \draw[->,thick,layerjade] (a3) -- (bc);
    \node[below] at (0.95,-1.75) {bijective};
  \end{scope}
  \begin{scope}[xshift=3.4cm]
    \foreach \i/\y in {1/1.2,2/0.4,3/-0.4} {
      \node[kotak] (c\i) at (0,\y) {\i};
    }
    \node[kotak] (da) at (1.9,1.2) {$a$};
    \node[kotak] (db) at (1.9,0.4) {$b$};
    \node[kotak] (dc) at (1.9,-0.4) {$c$};
    \node[kotak] (dd) at (1.9,-1.2) {$d$};
    \draw[->,thick,layerindigo] (c1) -- (da);
    \draw[->,thick,layerindigo] (c2) -- (db);
    \draw[->,thick,layerindigo] (c3) -- (dc);
    \node[below] at (0.95,-1.75) {injective only};
  \end{scope}
  \begin{scope}[xshift=6.9cm]
    \foreach \i/\y in {1/1.2,2/0,3/-1.2} {
      \node[kotak] (e\i) at (0,\y) {\i};
    }
    \node[kotak] (fa) at (1.9,0.8) {$a$};
    \node[kotak] (fb) at (1.9,-0.8) {$b$};
    \draw[->,thick,layerbrown] (e1) -- (fa);
    \draw[->,thick,layerbrown] (e2) -- (fa);
    \draw[->,thick,layerbrown] (e3) -- (fb);
    \node[below] at (0.95,-1.75) {surjective only};
  \end{scope}
\end{tikzpicture}
\caption{Three mappings between two columns of nodes. Left: every node is used exactly
once, hence bijective. Middle: the right column has four members and the node $d$
receives no incoming arrow, hence injective but not surjective. Right: two left nodes
go to $a$, hence surjective but not injective.}
\label{fig:bijeksi}
\end{figure}

The word "bijective" is only a short name for the combination of the two previous
properties. Its power appears in the following proposition: a bijection is the exact
condition for a function to be invertible.

\begin{proposisi}
$f : A \to B$ is bijective if and only if there is $g : B \to A$ with
$g \circ f = \mathrm{id}_A$ and $f \circ g = \mathrm{id}_B$. If $g$ exists, it is
unique.
\end{proposisi}
\par\smallskip\noindent\textit{Proof.} Suppose $g$ satisfies both equalities. If
$f(x) = f(y)$ then $x = g(f(x)) = g(f(y)) = y$, so $f$ is injective. For any
$b \in B$ we have $b = f(g(b))$, so $f$ is surjective.

Conversely, suppose $f$ is bijective. For each $b \in B$ there is exactly one
$x \in A$ with $f(x) = b$: existence from surjectivity, uniqueness from injectivity.
Define $g(b) = x$ for that $x$, then $g(f(x)) = x$ and $f(g(b)) = b$.

For uniqueness, suppose $g$ and $h$ satisfy both equalities. Then
\[
  g = g \circ \mathrm{id}_B = g \circ (f \circ h) = (g \circ f) \circ h
  = \mathrm{id}_A \circ h = h,
\]
which uses only the associativity of composition and the two equalities.
\hfill$\blacksquare${\parfillskip=0pt\par}

Composition behaves well with respect to all three properties: if $f$ and $g$ are
injective then $g \circ f$ is injective, and similarly for surjective and bijective.
The converse fails: with $f : \{1\} \to \{1,2\}$, $f(1) = 1$, and
$g : \{1,2\} \to \{1\}$, $g(1) = g(2) = 1$, the composition $g \circ f$ is bijective,
even though $g$ is not injective and $f$ is not surjective.

On finite sets the three properties lock together, and that is why "one-to-one" and
"onto" can be distinguished only on infinite sets.

\begin{lemma}
Let $A$ and $B$ be finite, $f : A \to B$, and $C \subseteq B$. Then
$|f(A)| \le |A|$ with equality if and only if $f$ is injective, and
$|C| \le |B|$ with equality if and only if $C = B$.
\end{lemma}
\par\smallskip\noindent\textit{Proof.} We prove the second part by induction on
$n = |B|$. For $n = 0$ we have $B = \varnothing$, hence $C = B$. Assume the claim for
sets with $n$ members and suppose $|B| = n+1$; choose $b \in B$ and write
$B' = B \setminus \{b\}$.

If $b \notin C$, then $C \subseteq B'$ and the induction hypothesis gives
$|C| \le |B'| = n < n+1$, with equality impossible. If $b \in C$, write
$C' = C \setminus \{b\} \subseteq B'$. Then $|C| = |C'| + 1 \le |B'| + 1 = |B|$;
equality holds exactly when $|C'| = |B'|$, which by the induction hypothesis means
$C' = B'$, that is $C = B$.

The first part is obtained by the same induction on $|A|$: write
$A = A' \cup \{a\}$ with $a \notin A'$ and observe that
$f(A) = f(A') \cup \{f(a)\}$, so $|f(A)| \le |A'| + 1 = |A|$. Equality demands
$f(a) \notin f(A')$ and $|f(A')| = |A'|$; the latter means $f|_{A'}$ is injective by
the induction hypothesis, and the former means $f$ does not repeat a value at $a$.
\hfill$\blacksquare${\parfillskip=0pt\par}

\begin{proposisi}
Let $A$ and $B$ be finite with $|A| = |B|$, and let $f : A \to B$. Then $f$
injective, $f$ surjective, and $f$ bijective are three equivalent statements.
\end{proposisi}
\par\smallskip\noindent\textit{Proof.} If $f$ is injective, the Lemma gives
$|f(A)| = |A| = |B|$; since $f(A) \subseteq B$ and both have the same size, the Lemma
gives $f(A) = B$, so $f$ is surjective. If $f$ is surjective, then $f(A) = B$ so that
$|f(A)| = |B| = |A|$, and the Lemma forces $f$ to be injective. Since bijective means
both properties at once, the equivalence of the three is complete.
\hfill$\blacksquare${\parfillskip=0pt\par}

A warning to carry into the next chapter: that proposition uses the word "finite"
essentially. On $\mathbb{N}$ the function $f(x) = 2x$ is injective but not
surjective, and the function mapping $1$ to $1$ and $n \ge 2$ to $n-1$ is surjective
but not injective.

\section{Finite, countable, and uncountable sets}
\tingkat{advanced}

Now we return to the opening question. To compare two sets without counting we pair
their members, and all we need to require is that the pairing uses an invertible
function.

\begin{definisi}[Equinumerous]
The sets $A$ and $B$ are equinumerous, written $A \sim B$, if there is a bijection
$A \to B$.
\end{definisi}

The relation $\sim$ is reflexive (the identity is a bijection), symmetric (the inverse
of a bijection is a bijection), and transitive (the composition of two bijections is a
bijection). So it behaves as an equivalence relation, but not on "the set of all sets",
because that collection is not a set; this is discussed in the box for researchers.

\begin{definisi}[Finite, countable, uncountable]\index{countable set}
A set $A$ is finite if $A = \varnothing$ or there is $n \in \mathbb{N}$ and a
bijection $A \to \{1,2,\dots,n\}$; in the second case we write $|A| = n$. A set $A$
is countable if $A \sim \mathbb{N}$. An infinite set that is not countable is called
uncountable.
\end{definisi}

\kamus{countable}{terhitung}{All its members can be given serial numbers with none skipped, as when every passenger receives a distinct ticket number.}

The name "countable" is apt: if $A \sim \mathbb{N}$, the bijection is a list. We can
name the first member, the second, and so on, and every member appears exactly once.
What is surprising is how many sets can be listed that way. Pairs of numbers, for
instance.

\begin{proposisi}[Cantor pairing]\index{Cantor pairing}
$\mathbb{N} \times \mathbb{N}$ is countable.
\end{proposisi}
\par\smallskip\noindent\textit{Proof.} Define
$\pi(m,n) = \frac{(m+n)(m+n+1)}{2} + n$ for $m, n \in \mathbb{N} =
\{0,1,2,\dots\}$. Write $T_s = \frac{s(s+1)}{2}$, the number of pairs with
$m + n < s$. For a pair with $m + n = s$ we have $0 \le n \le s$, so
$\pi(m,n) = T_s + n$ takes the values $T_s, \dots, T_s + s$ exactly once each. Since
$T_s + s = T_{s+1} - 1$, the block of values for diagonal $s$ is
$\{T_s, \dots, T_{s+1}-1\}$. The blocks join without gaps and without overlap
starting from $T_0 = 0$. Hence every counting number lies in exactly one block, and
within its block comes from exactly one pair; therefore $\pi$ is bijective.
\hfill$\blacksquare${\parfillskip=0pt\par}

The inverse is explicit: $s = \lfloor (\sqrt{8q+1}-1)/2 \rfloor$, $n = q - T_s$, and
$m = s - n$. That square root is the one place where machine arithmetic can deceive
us, and the example is in the lab section.

Figure~\ref{fig:kisi} draws that grid: the cell in row $m$ and column $n$ is labelled
$\pi(m,n)$, and the color marks the diagonal $m + n = s$, so that reading by color
gives $0$ through $14$ without gaps and without repetition. The right panel applies
the same reading order to a list of four real numbers, with row $d$ flipping the
diagonal digits so that $d$ is not in the list.

\begin{figure}[htbp]
\centering
\begin{tikzpicture}[scale=0.95,every node/.style={font=\footnotesize}]
  \foreach \s/\warna in {0/layerteal,1/layerorange,2/layerindigo,3/layerjade,4/layerplum} {
    \foreach \n in {0,...,\s} {
      \pgfmathtruncatemacro{\m}{\s-\n}
      \pgfmathtruncatemacro{\v}{int(\s*(\s+1)/2+\n)}
      \draw[fill=\warna!30,draw=ink!45] (\n*0.86,-\m*0.72) rectangle ++(0.8,0.66);
      \node at (\n*0.86+0.4,-\m*0.72+0.33) {\v};
    }
  }
  \draw[->,thick] (4.35,0.33) -- (4.95,0.33);
  \node[right] at (4.95,0.33) {$n$};
  \draw[->,thick] (-0.45,-0.05) -- (-0.45,-3.0) node[below] {$m$};
  \begin{scope}[xshift=6.6cm,yshift=0.45cm]
    \foreach \i/\j/\b in {1/1/0,1/2/1,1/3/1,1/4/0, 2/1/1,2/2/0,2/3/0,2/4/1,
                          3/1/0,3/2/0,3/3/1,3/4/0, 4/1/1,4/2/1,4/3/0,4/4/0} {
      \ifnum\i=\j
        \fill[accent!28] (\j*0.52-0.52,-\i*0.42) rectangle ++(0.5,0.38);
      \fi
      \draw[ink!45] (\j*0.52-0.52,-\i*0.42) rectangle ++(0.5,0.38);
    }
    \draw[accent,thick] (0.25,-0.23) -- (1.79,-1.49);
    \foreach \i/\j/\b in {1/1/0,1/2/1,1/3/1,1/4/0, 2/1/1,2/2/0,2/3/0,2/4/1,
                          3/1/0,3/2/0,3/3/1,3/4/0, 4/1/1,4/2/1,4/3/0,4/4/0} {
      \node at (\j*0.52-0.27,-\i*0.42+0.19) {\b};
    }
    \foreach \i in {1,...,4} {
      \node[left] at (-0.62,-\i*0.42+0.19) {$x_{\i}$};
    }
    \foreach \j/\b in {1/1,2/1,3/0,4/1} {
      \fill[layergrey!25] (\j*0.52-0.52,-2.34) rectangle ++(0.5,0.38);
      \draw[ink!45] (\j*0.52-0.52,-2.34) rectangle ++(0.5,0.38);
      \node at (\j*0.52-0.27,-2.34+0.19) {\b};
    }
    \node[left] at (-0.10,-2.15) {$d$};
    \node[below] at (0.52,-2.42) {digit flipped};
  \end{scope}
\end{tikzpicture}
\caption{Left: the grid of pairs $(m,n)$ with $m + n \le 4$, each cell labelled
$\pi(m,n)$, with color marking the diagonal $m + n = s$; the first five diagonals
carry the values $0$ through $14$. Right: the digit table of four real numbers, with
the diagonal cells colored and row $d$ flipping the diagonal digits $0,0,1,0$ into
$1,1,0,1$.}
\label{fig:kisi}
\end{figure}

Since every positive rational is a pair of natural numbers already in lowest terms,
the set of rationals is countable too. All we must guard against is listing a number
twice.

\begin{proposisi}
The set $\mathbb{Q}_{>0}$ of positive rational numbers is countable.
\end{proposisi}
\par\smallskip\noindent\textit{Proof.} Every $q \in \mathbb{Q}_{>0}$ can be written
uniquely as $q = p/r$ with $p, r$ natural numbers and $\gcd(p,r) = 1$. Define the
height of $q$ as $p + r$. For each $k \ge 2$ there are only finitely many rationals of
height $k$, since $p$ and $r$ lie between $1$ and $k-1$. The following list therefore
enumerates $\mathbb{Q}_{>0}$ without repetition: for $k = 2, 3, 4, \dots$ write down
all $p/r$ of height $k$ with $\gcd(p,r) = 1$, increasing in $p$. Every positive
rational appears exactly once, namely at level $k = p + r$ in its unique reduced form,
so that list is a bijection $\mathbb{N} \to \mathbb{Q}_{>0}$.
\hfill$\blacksquare${\parfillskip=0pt\par}

\begin{akibat}
$\mathbb{Q}$ and $\mathbb{Z}$ are countable.
\end{akibat}
\par\smallskip\noindent\textit{Proof.} The list $0, 1, -1, 2, -2, \dots$ is a
bijection $\mathbb{N} \to \mathbb{Z}$. For $\mathbb{Q}$, enumerate each nonzero
rational as $\pm p/r$ in reduced form by $(p+r)$ with alternating signs, then place
$0$ at the front; one extra member does not destroy countability.
\hfill$\blacksquare${\parfillskip=0pt\par}

So $\mathbb{N}$, $\mathbb{Z}$, and $\mathbb{Q}$ are equinumerous, even though
$\mathbb{Q}$ feels far more crowded. Density is not a measure of size; the measure is
whether or not a list exists.

\begin{tantangan}
A hotel on the seafront has rooms numbered $1, 2, 3, \dots$ without end, and that
night every room holds exactly one guest. A traveler arrives without a reservation.
The receptionist says the hotel is full; the manager says there is still a room for
him without sending a single guest home.

Show a rule for moving guests that keeps the manager's promise, then explain why that
rule is a bijection between the old guests plus one new guest and all the rooms. Why
can the same hotel not accommodate as many new guests as there are real numbers?

Hint 1: do not look for an empty room, because there is none. Create an empty room by
moving its occupant.

Hint 2: shift the occupant of room $n$ to room $n+1$ for every $n$, then put the
traveler in room $1$; check that every room holds exactly one person and that no
guest is lost.

If you get stuck, return to the section Finite, countable, and uncountable sets, in
particular the definition of equinumerous.
\end{tantangan}

Cantor showed that $\mathcal{P}(A)$ is always larger than $A$, by a trick that does not
depend on the size of $A$.

\begin{teorema}[Cantor]\index{Cantor diagonal}
For every set $A$ there is no surjective function $A \to \mathcal{P}(A)$.
\end{teorema}
\par\smallskip\noindent\textit{Proof.} Let $f : A \to \mathcal{P}(A)$ be arbitrary
and define $D = \{\, a \in A : a \notin f(a) \,\}$. Then $D \subseteq A$, so
$D \in \mathcal{P}(A)$. Suppose $D = f(a_0)$ for some $a_0 \in A$. If
$a_0 \in D$, then by the definition of $D$ we have $a_0 \notin f(a_0) = D$, a
contradiction. If $a_0 \notin D$, then $a_0 \notin f(a_0)$, which is precisely the
membership condition for $D$, so $a_0 \in D$, a contradiction again. Hence $D$ is not
in the range of $f$.
\hfill$\blacksquare${\parfillskip=0pt\par}

The set $D$ is formed by asking, for each $a$, whether $a$ is a member of the set
paired with it, then flipping the answer; that flip is what makes $D$ differ from
every value of $f$. For $A = \mathbb{N}$ the trick appears as reading the diagonal of
a table with every entry changed, which is why it is called the diagonal argument.

\begin{akibat}
$\mathcal{P}(\mathbb{N})$ is uncountable, and $\mathbb{R}$ is uncountable.
\end{akibat}
\par\smallskip\noindent\textit{Proof.} By Cantor's Theorem with $A = \mathbb{N}$,
there is no surjection $\mathbb{N} \to \mathcal{P}(\mathbb{N})$, so
$\mathcal{P}(\mathbb{N})$ is uncountable.

For $\mathbb{R}$ it suffices to show that $[0,1]$ is uncountable. Every $x \in [0,1]$
has a binary expansion $x = \sum_{k \ge 1} b_k 2^{-k}$ with $b_k \in \{0,1\}$; its
existence is proved in Chapter 3 and Chapter 5. The map from the set of binary
sequences to $[0,1]$ is therefore surjective, and the set of binary sequences is
equinumerous with $\mathcal{P}(\mathbb{N})$ because a binary sequence is a function
$\mathbb{N} \to \{0,1\}$. If $[0,1]$ were countable, its preimage in the set of binary
sequences would be countable as well, contradicting the previous paragraph.
\hfill$\blacksquare${\parfillskip=0pt\par}

One note of honesty: binary expansions are not unique, since
$0.0111\dots = 0.1000\dots$ in base two. The argument above is undisturbed because
it uses only the surjectivity, but $|\mathbb{R}| = |\mathcal{P}(\mathbb{N})|$ demands
the following tool.

\begin{proposisi}[Cantor--Bernstein--Schroeder]
If there is an injective function $A \to B$ and an injective function $B \to A$, then
$A \sim B$.
\end{proposisi}
\par\smallskip\noindent\textit{Proof sketch.} Let $f : A \to B$ and $g : B \to A$ be
injective. Let $A_0$ be the set of points of $A$ that have no predecessor in $B$, then
$A_{k+1} = g(f(A_k))$, and write $A^{+}$ for the union of all $A_k$ with $k$ even.
Define $h(x) = f(x)$ if $x \in A^{+}$ and $h(x) = g^{-1}(x)$ if $x \notin A^{+}$. On
the first part $h$ is injective because $f$ is injective, on the second because
$g^{-1}$ is injective, and the images of the two parts are separated by the way
$A^{+}$ is built, so $h$ is injective. For surjectivity of the image, take
$y \in B$. If $g(y) \in A^{+}$, then $g(y)$ lies in some $A_k$ with $k$ even so that
$g(y) = g(f(x))$ for $x \in A_{k-1}$ and $h(x) = y$; if $g(y) \notin A^{+}$, then
$h(g(y)) = y$. Hence $h$ is surjective.

The bookkeeping of chain depth is the delicate part, and its details are skipped here;
the full formal version is in \emph{Mathematics in Lean}, section 4.3.
\hfill$\blacksquare${\parfillskip=0pt\par}

With that tool we may conclude $|\mathbb{R}| = |\mathcal{P}(\mathbb{N})|$: an
injection from $\mathcal{P}(\mathbb{N})$ to $\mathbb{R}$ is given by binary expansion,
and the injection the other way by the tangent function. The size of $\mathbb{R}$ is
usually written $2^{\aleph_0}$, with $\aleph_0$ the size of $\mathbb{N}$.

\begin{untukanak}
Picture the ticket machine at the entrance to an exhibition. Every visitor receives one
numbered ticket: 1, 2, 3, and so on, no number is used twice, and every possible
number will eventually be given out. A collection of things that can be ticketed like
that is called countable.

It looks as though every endless collection can be ticketed. But not so. Some
collections have so many members that however many tickets we print, there is always a
member left without one. The ticket machine is not broken; that collection really is
bigger than the supply of numbers.

The question that makes one curious: how can we know that a collection cannot be
ticketed, when not all methods have been tried? The answer is Cantor's flipping trick.
Each time someone claims to have ticketed all the members, we build one new member
deliberately different from every ticketed member; that member is certainly left out.
\end{untukanak}

\section{Inequalities, absolute value, and upper bounds}
\tingkat{intermediate}

The last section installs two tools for Chapter 3. The first measures the distance
between two numbers regardless of direction, namely the absolute value.

\begin{definisi}[Absolute value]\index{absolute value}
For $x \in \mathbb{R}$, $|x| = x$ if $x \ge 0$ and $|x| = -x$ if $x < 0$.
\end{definisi}

A route through a third point is never shorter than the direct route.

\begin{proposisi}[Triangle inequality]
For every $x, y \in \mathbb{R}$, $|x + y| \le |x| + |y|$ and
$\bigl| |x| - |y| \bigr| \le |x - y|$.
\end{proposisi}
\par\smallskip\noindent\textit{Proof.} Note that $-|x| \le x \le |x|$ and
$-|y| \le y \le |y|$. Adding the two gives
$-(|x|+|y|) \le x + y \le |x|+|y|$, which by the definition of absolute value means
$|x+y| \le |x|+|y|$. For the second, write $x = (x-y) + y$ and apply the first result:
$|x| \le |x-y| + |y|$, so $|x| - |y| \le |x-y|$. Exchanging the roles of $x$ and $y$
gives $|y| - |x| \le |x-y|$, hence $\bigl| |x| - |y| \bigr| \le |x-y|$.
\hfill$\blacksquare${\parfillskip=0pt\par}

Now bounds. A set of numbers can be bounded above without having a largest member, and
precisely that case is the most important one.

\begin{definisi}[Upper bound and supremum]\index{supremum}
Let $S \subseteq \mathbb{R}$ be nonempty. A number $u$ is called an upper bound of $S$
if $s \le u$ for every $s \in S$; if one exists, $S$ is called bounded above. A number
$u$ is called the supremum of $S$, written $u = \sup S$, if $u$ is an upper bound of
$S$ and $u \le v$ for every upper bound $v$ of $S$.
\end{definisi}

\kamus{supremum}{batas atas terkecil}{The tightest fence: the highest boundary that cannot be lowered any further without leaving a member above it.}

The supremum, when it exists, is unique, since two least upper bounds dominate each
other. What the definition does not guarantee is existence:
$\{x \in \mathbb{Q} : x^{2} < 2\}$ is bounded above in $\mathbb{Q}$ but has no
supremum there.

\begin{catatan}[Completeness axiom]
Every nonempty set $S \subseteq \mathbb{R}$ that is bounded above has a supremum in
$\mathbb{R}$. This axiom is not proved here; its construction together with its
consequences is the content of Chapter 3, including the epsilon characterization that
turns a statement about a least bound into a statement about $\varepsilon$, and the
example of a set that has a supremum without a maximum.
\end{catatan}

On a finite set the supremum and the maximum coincide, and swapping the two often
slips through on small examples.

\section{Worked examples}
\tingkat{intermediate}

The following four computations are done by hand down to the number, then checked
again with a script. All of them are small, so arithmetic errors show up immediately.

\begin{resep}
For two finite sets $A$ and $B$:
\begin{enumerate}
  \item Compute $|A|$, the number of members of $A$.
  \item Compute $|B|$ the same way.
  \item Compute $|A \cap B|$, the members appearing in both lists.
  \item Add and then subtract: $|A \cup B| = |A| + |B| - |A \cap B|$.
\end{enumerate}
A warning about order: step 3 may not be postponed until after step 4. If the addition
is done first, the members of the intersection are counted twice and the result
overshoots by exactly $|A \cap B|$. For $A = \{1,\dots,5\}$ and $B = \{4,\dots,8\}$:
$5 + 5 - 2 = 8$, the same number as in the Venn diagram above.
\end{resep}

First, the power set. For $A = \{1,2,3\}$ the subsets have sizes $0,1,2,3$ with counts
$1, 3, 3, 1$, in total $1 + 3 + 3 + 1 = 8 = 2^{3}$. The row $n = 4$ gives
$1+4+6+4+1 = 16$, $n = 2$ gives $1+2+1 = 4$, and $n = 1$ gives $1+1 = 2$: the same
Pascal triangle as in the Proposition about $2^{n}$.

Second, the Cantor pairing. With $\pi(m,n) = T_{m+n} + n$ and $T_0 = 0$, $T_1 = 1$,
$T_2 = 3$, $T_3 = 6$: diagonal $s = 0$ contains $(0,0) \mapsto 0$; $s = 1$ contains
$(1,0) \mapsto 1$ and $(0,1) \mapsto 2$; $s = 2$ contains $(2,0) \mapsto 3$,
$(1,1) \mapsto 4$, $(0,2) \mapsto 5$; $s = 3$ contains $(3,0) \mapsto 6$,
$(2,1) \mapsto 7$, $(1,2) \mapsto 8$, $(0,3) \mapsto 9$. Those ten values are exactly
$0, \dots, 9$ without repetition and without gaps.

The opposite direction is easier to get wrong. Take $q = 1000$. We look for $s$ with
$T_s \le 1000 < T_{s+1}$. Since $T_{44} = 44 \cdot 45/2 = 990$ and $T_{45} = 1035$, we
get $s = 44$, then $n = 1000 - 990 = 10$ and $m = 44 - 10 = 34$. Check:
$\pi(34,10) = 44 \cdot 45/2 + 10 = 990 + 10 = 1000$, correct.

Third, the rational numbers. By the height rule, level $k = 2$ contains $1/1$;
$k = 3$ contains $1/2, 2/1$; $k = 5$ contains $1/4, 2/3, 3/2, 4/1$. The number of terms
at each level is the value of Euler's totient function, $1, 2, 2, 4, 2, 6$ for
$k = 2, \dots, 7$. The cumulative count up to $k = 7$ is $1+2+2+4+2+6 = 17$, so $6/1$
occupies position 17. Since the cumulative count up to $k = 9$ is $27$ and level
$k = 10$ contains $1/9, 3/7, 7/3, 9/1$, the number $7/3$ occupies position 30.

Fourth, the number of relations. On a set with $3$ members there are $2^{9} = 512$
relations, and only $5$ are equivalence relations: $1$ with one class, $3$ with two
classes, $1$ with three classes. On a set with $4$ members there are $2^{16} = 65{,}536$
relations and $15$ that are equivalence relations. Those four numbers appear again as
the last four lines of the script output; the sequence $1, 2, 5, 15, 52$ is known as
the Bell numbers.

\kertas{Four numbers already checked again. (i) $|\mathcal{P}(A)| = 2^{n}$ for
$n = 0,1,2,3,4$ gives $1, 2, 4, 8, 16$.
(ii) $\pi(34,10) = 990 + 10 = 1000$, and $\pi$ on diagonals $s = 0,1,2,3$ gives
$0,1,2,3,4,5,6,7,8,9$. (iii) The number of positive rationals of height $k$ is the
totient value at $k$; the cumulative count up to $k = 9$ is $27$, so $7/3$ sits at
position 30. (iv) Relations on a set with $3$ members: $512$ in all, $5$ equivalence
relations.}

\section{Computational lab}
\tingkat{advanced}

The following script runs the four computations and adds one check that cannot be done
by hand: it compares the inverse of the Cantor pairing computed with machine floating
point against the one computed with integer square roots. Both should give the same
answer; at one point the answers differ.

\begin{lab}
\begin{verbatim}
from itertools import combinations
from math import gcd, isqrt, sqrt, floor

# (1) number of subsets for small n
for n in range(5):
    A = range(n)
    subs = [S for k in range(n + 1) for S in combinations(A, k)]
    print("n =", n, "|P(A)| =", len(subs), "2^n =", 2 ** n)

# (2) Cantor pairing pi(m, n) = (m+n)(m+n+1)/2 + n
def pi(m, n):
    s = m + n
    return s * (s + 1) // 2 + n

pasangan = [(0, 0), (1, 0), (0, 1), (2, 0), (1, 1),
            (0, 2), (3, 0), (2, 1), (1, 2), (0, 3)]
print("pi:", [pi(m, n) for m, n in pasangan])
keys = {pi(m, n) for m in range(61) for n in range(61)}
print("unique keys over 61x61 pairs:", len(keys) == 61 * 61)

# (3) inverse: float vs isqrt
s, n = 2 * 10**8, 2 * 10**8 - 1
q = s * (s + 1) // 2 + n
print("q =", q, "s correct =", s)
print("s from float:", int(floor((sqrt(8 * q + 1) - 1) / 2)))
print("s from isqrt:", (isqrt(8 * q + 1) - 1) // 2)

# (4) enumerate positive Q by height k = p + q
seq = []
for k in range(2, 10):
    row = [(p, k - p) for p in range(1, k) if gcd(p, k - p) == 1]
    print("k =", k, "count =", len(row), row)
    seq += row
print("terms 1..8:", seq[:8])
print("total up to k = 9:", len(seq))

# (5) relations on an n-element set: 2^(n*n) in all, Bell numbers equivalence
def ekuivalen(R, n):
    refl = all((a, a) in R for a in range(n))
    sim = all((b, a) in R for a, b in R)
    trans = all((a, c) in R for a, b in R for c in range(n)
                if (b, c) in R)
    return refl and sim and trans

for n in range(1, 5):
    semua = 2 ** (n * n)
    hitung = 0
    for bit in range(semua):
        R = {(a, b) for a in range(n) for b in range(n)
             if bit >> (a * n + b) & 1}
        if ekuivalen(R, n):
            hitung += 1
    print("n =", n, "relations =", semua, "equivalences =", hitung)
\end{verbatim}
\end{lab}

The output as printed, verbatim:

\begin{verbatim}
n = 0 |P(A)| = 1 2^n = 1
n = 1 |P(A)| = 2 2^n = 2
n = 2 |P(A)| = 4 2^n = 4
n = 3 |P(A)| = 8 2^n = 8
n = 4 |P(A)| = 16 2^n = 16
pi: [0, 1, 2, 3, 4, 5, 6, 7, 8, 9]
unique keys over 61x61 pairs: True
q = 20000000299999999 s correct = 200000000
s from float: 200000001
s from isqrt: 200000000
k = 2 count = 1 [(1, 1)]
k = 3 count = 2 [(1, 2), (2, 1)]
k = 4 count = 2 [(1, 3), (3, 1)]
k = 5 count = 4 [(1, 4), (2, 3), (3, 2), (4, 1)]
k = 6 count = 2 [(1, 5), (5, 1)]
k = 7 count = 6 [(1, 6), (2, 5), (3, 4), (4, 3), (5, 2), (6, 1)]
k = 8 count = 4 [(1, 7), (3, 5), (5, 3), (7, 1)]
k = 9 count = 6 [(1, 8), (2, 7), (4, 5), (5, 4), (7, 2), (8, 1)]
terms 1..8: [(1, 1), (1, 2), (2, 1), (1, 3), (3, 1), (1, 4), (2, 3), (3, 2)]
total up to k = 9: 27
n = 1 relations = 2 equivalences = 1
n = 2 relations = 16 equivalences = 2
n = 3 relations = 512 equivalences = 5
n = 4 relations = 65536 equivalences = 15
\end{verbatim}

What must be read from that output: the first five lines confirm $2^{n}$ for
$n = 0$ through $4$; the ten values of $\pi$ are consecutive without gaps, and
$61 \times 61 = 3721$ pairs give $3721$ unique keys, which is numerical evidence, not a
mathematical proof, because the mathematical proof is the $T_s$ block argument. The
last four lines sweep all $2^{n^{2}}$ relations for $n = 1$ through $4$ and count those
that are reflexive, symmetric, and transitive: the result is $1, 2, 5, 15$. There is no
line for $n = 5$ because sweeping $2^{25}$ relations takes too long; its value is $52$.

Third, the inverse with faulty floating point. At $q = 20{,}000{,}000{,}299{,}999{,}999$ the
formula $\lfloor (\sqrt{8q+1}-1)/2 \rfloor$ with 64-bit machine floating point gives
$s = 200{,}000{,}001$, whereas the true value is $200{,}000{,}000$. The error is one, and that
is enough to return a completely different pair: with $s = 200{,}000{,}001$ the value
$n = q - T_s$ becomes negative. The number $8q+1$ is of order $10^{17}$, whereas
64-bit floating point stores only about $16$ decimal digits exactly; \texttt{isqrt}
works on integers and gives $200{,}000{,}000$. The lesson is general: inverting an exact
numbering formula demands integer arithmetic, as discussed in Chapter 22.

\begin{untukdosen}
Four meetings. First: set operations and De Morgan, with each proof written as a chain
of equivalences, not as two separate inclusions. Second: Cartesian products,
relations, then the filtering into partial orders and equivalence relations. Third:
classes, partitions, and functions, with the Lemma on classes proved at the board.
Fourth: cardinality up to Cantor's Theorem, plus a sketch of Cantor--Bernstein if time
allows.

Two misconceptions appear most often. First, students think a symmetric and transitive
relation is automatically reflexive; the counterexample
$R = \{(1,1),(2,2),(1,2),(2,1)\}$ on $\{1,2,3\}$ is usually enough, and Exercise 2 asks
for a verdict on that claim. Second, students think $\mathbb{Q}$ is larger than
$\mathbb{N}$ because of its density; ask for the first ten terms of the height
enumeration at the board, then ask for the position of $7/3$.

An activity that works: divide the board into "natural numbers" and "positive
rationals", then have two volunteers pair members for two minutes. The naive pairing
will fail, and then together build the pairing by height. Students are ready to move on
when they can explain why that list misses no rational, and why a similar list fails
for $[0,1]$.
\end{untukdosen}

\section{Summary}
\tingkat{basic}

\begin{ringkasan}
\begin{itemize}
  \item Union, intersection, difference, and complement are defined through
        membership conditions; De Morgan's laws are laws of logic transferred to
        sets, and their version holds for any indexed family.
  \item $|\mathcal{P}(A)| = 2^{|A|}$ for finite $A$, via a bijection to the
        characteristic functions $A \to \{0,1\}$.
  \item A relation is a subset of $A \times B$; on a set with $n$ members there are
        $2^{n^{2}}$ relations, and four properties filter them into equivalence
        relations and partial orders.
  \item Equivalence relations and partitions are two faces of the same object; the
        classes are always nonempty, mutually disjoint, and cover the whole set.
  \item A bijective function is exactly one that has a unique two-sided inverse, and
        on finite sets of equal size the three properties lock together.
  \item $\mathbb{N} \times \mathbb{N}$ and $\mathbb{Q}$ are countable through explicit
        numberings, whereas $\mathcal{P}(\mathbb{N})$ and $\mathbb{R}$ are uncountable
        through Cantor's diagonal argument.
  \item The supremum is the least upper bound; its existence is guaranteed by the
        completeness axiom, while the epsilon characterization and its construction
        continue in Chapter 3.
\end{itemize}
\end{ringkasan}

\section{Exercises}
\tingkat{intermediate}

\begin{latihan}
\begin{enumerate}
  \item Prove $A \cap (B \cup C) = (A \cap B) \cup (A \cap C)$ and
        $A \setminus (B \setminus C) = (A \setminus B) \cup (A \cap C)$ with chains of
        membership equivalences.
  \item Assess this claim: true, false, or not enough evidence. "Every relation on a
        nonempty set that is symmetric and transitive is necessarily reflexive." If it
        is false, give the smallest possible counterexample; if there is not enough
        evidence, state what further evidence would be needed.
  \item Prove that if $f : A \to B$ is injective, then $S \mapsto f(S)$ is an
        injective function from $\mathcal{P}(A)$ to $\mathcal{P}(B)$. Explain where
        injectivity is used and why the conclusion fails without it.
  \item Prove Kuratowski's ordered pair lemma:
        $\{\{a\},\{a,b\}\} = \{\{c\},\{c,d\}\}$ if and only if $a = c$ and
        $b = d$. Handle the case $a = b$ separately.
  \item Complete the first part of the counting Lemma: for $f : A \to B$ with $A$
        finite, $|f(A)| \le |A|$ and equality holds only if $f$ is injective.
  \item Count the equivalence relations on a set with $5$ members that have exactly
        two classes, then check the answer with Python.
  \item Compute $\pi(7,5)$ with the Cantor pairing formula, then find $(m,n)$ with
        $\pi(m,n) = 1000$. Check both answers with Python.
  \item Determine the position of $7/3$ in the enumeration of positive rationals by
        height, then determine the number at positions 30 and 41. Check with Python.
  \item Prove that if $A_1, A_2, A_3, \dots$ are each countable and the numbering of
        every $A_n$ is given explicitly, then $\bigcup_{n \ge 1} A_n$ is countable.
  \item Show that $(0,1)$ is equinumerous with $\mathbb{R}$ through one explicit
        bijection, and explain why the function we choose must have its domain
        restricted.
\end{enumerate}
\end{latihan}

Hints for the three hardest items. Exercise 5: in the induction step write
$A = A' \cup \{a\}$ with $a \notin A'$, then compare $f(A)$ with
$f(A') \cup \{f(a)\}$; equality of sizes demands $f(a) \notin f(A')$ and
$|f(A')| = |A'|$ at the same time. Exercise 6: the number of ways to split $n$ members
into two nonempty parts is $(2^{n}-2)/2$, since choosing the first part or the second
gives the same split. Exercise 9: arrange the numbering of all the $A_n$ as a two-column
table, read it diagonal by diagonal as in the proof that $\mathbb{N} \times \mathbb{N}$
is countable, then discard repeated terms.

\begin{untukpeneliti}
Four things are genuinely proved here. Cantor's Theorem uses no axiom of choice and is
constructive: from any $f$ we build $D$ that lies outside its range. The
uncountability of $\mathbb{R}$ depends on the binary expansion of every point of
$[0,1]$; that fact requires completeness and here is only referred to Chapter 3 and
Chapter 5.

The Cantor--Bernstein--Schroeder theorem in this chapter is only a sketch. A full
proof requires careful bookkeeping of chain depth; the complete formal version is in
Avigad--Massot, \emph{Mathematics in Lean}, section 4.3, while the classical treatment
of sets, functions, and equivalence relations is in Hoffman--Kunze,
\emph{Linear Algebra}, Appendix A.1--A.3. The primary references for the diagonal and
for uncountability are Cantor, "Ueber eine Eigenschaft des Inbegriffes aller
reellen algebraischen Zahlen" (1874) and "Ueber eine elementare Frage der
Mannigfaltigkeitslehre" (1891).

The most famous open question: is there a cardinal between $|\mathbb{N}|$ and
$|\mathbb{R}|$? The continuum hypothesis states that there is none. Godel (1940) showed
that the hypothesis is consistent with ZFC, Cohen (1963) showed that its negation is
consistent too. So the question cannot be answered inside ZFC.

Two technical limits. First, $\sim$ may not be called an equivalence relation on "the
set of all sets", because that collection is not a set; the correct treatment uses
classes as metatheoretic objects. Second, the corollary "a countable union of
countable sets is countable" requires countable choice when the numberings are not
given, and Exercise 9 deliberately requires an explicit numbering.
\end{untukpeneliti}

\peta{Chapter 2 maps to Abbott, \emph{Understanding Analysis}, sections "Some
Preliminaries" (sets, functions, cardinality) and "Cantor's Theorem"; Hoffman--Kunze,
\emph{Linear Algebra}, Appendix A.1 (sets), A.2 (functions and inverses), A.3
(equivalence relations, classes, quotients); Avigad--Massot, \emph{Mathematics in
Lean}, Chapter 4 ("Sets and Functions") for extensionality, preimages, and
Cantor--Bernstein--Schroeder, and Chapter 6.1 for finite sets.}
